%! Justifying a Claim Based on a Confidence Interval for a Difference of Population Proportions — Unit 6, Lesson 6.9 — Homework
\documentclass[10pt]{article}
\usepackage{apstats-article}
\usepackage{apstats-boxes}
\newcommand{\TallMath}[1]{$\displaystyle #1\rule[-1.4em]{0pt}{3.2em}$}

\begin{document}

\pageheader{Unit 6, Lesson 6.9}{Homework}
\namedateperiod

\vspace{0.15in}

\begin{enumerate}

\item \textbf{Interpret and Evaluate.} A researcher found that a 99\% CI for
$p_{\text{city}} - p_{\text{suburb}}$, the difference in the proportions of residents
who use public transit, is $(-0.06, 0.22)$.

\begin{enumerate}[label=\alph*.]
  \item Interpret this interval. \writelines{2}
  \item Does this interval provide evidence that city residents use public transit at a
        higher rate than suburban residents? Explain. \writelines{2}
  \item Would you reach the same conclusion with an 80\% CI? Why or why not?
        \writelines{2}
\end{enumerate}

\item \textbf{Full Procedure.} Independent random samples were taken from two
hospitals to compare the proportions of patients discharged within 3 days of admission.
Hospital A: $n_1=400$, $\hat p_1 = 0.73$. Hospital B: $n_2=350$, $\hat p_2 = 0.65$.

\begin{enumerate}[label=\alph*.]
  \item Verify all three conditions. \writelines{3}
  \item Construct a 95\% CI for $p_A - p_B$. \writelines{3}
  \item Interpret the interval and evaluate whether there is evidence of a difference
        in discharge rates. \writelines{2}
\end{enumerate}

\end{enumerate}

\end{document}
